\chapter{The Ledger}\label{ch:ledger}

Everything so far has been defined over a finite set of arguments, attacks and policies. This chapter says where that set comes from and what may be done to it. The answer is an append-only ledger of records, from which the argument graph is derived and from which the displayed dossier is derived in turn. Keeping the three apart is the central discipline:
\[
\text{ledger}\ \longrightarrow\ \text{argument graph}\ \longrightarrow\ \text{epistemic presentation}.
\]
The ledger records what happened. The argument graph represents inferential relations among what was admitted. The presentation is an interpretation of that graph under a stated policy. Many systems fail because a single interface label such as ``accepted'' is treated at once as a historical fact, a logical consequence and a community consensus. In Adversaria they are three different objects with three different sources of authority.

\section{Records and the pipeline}

Contributions pass through five distinct transitions:
\[
\textsc{Inscribe}\to\textsc{Validate}\to\textsc{Admit}\to\textsc{Project}\to\textsc{Activate}.
\]
\emph{Inscribe} stores the submission in the ledger. \emph{Validate} checks schema, integrity, references and anchoring. \emph{Admit} decides whether the record is a recognizable argumentative act. \emph{Project} decides whether it appears in a particular dossier. \emph{Activate} permits it to influence labels, status and search ranking. The first is the only one every submission receives without conditions. None entails the next:
\[
\mathrm{stored}(r)\ \not\Rightarrow\ \mathrm{admitted}(r)\ \not\Rightarrow\ \mathrm{endorsed}(r).
\]
This chapter formalizes the ledger, admission and activation. Projection and its governance belong to the protocol part; only what evaluation may depend on is fixed here (Section~\ref{sec:separation}).

\begin{definition}[Record]\label{def:record}
A \emph{record} is a tuple $r=(\mathsf{type},\mathsf{payload},\mathsf{refs},\mathsf{author},\mathsf{origin},\mathsf{trace})$. The type is one of claim, evidence, argument, policy, distinction, restriction, revision, equivalence, withdrawal, and so on. The payload carries the fields defined in Chapters~\ref{ch:claim}--\ref{ch:objection}. The references are identifiers of earlier records that the payload cites. The origin is one of $\{\text{human},\text{automated},\text{assisted},\text{imported}\}$. The trace is defined in Section~\ref{sec:trace}.
\end{definition}

\subsection{Three identifiers}

A record needs several notions of sameness, and they must not be conflated.

\begin{definition}[Identifiers]\label{def:ids}
Let $\Hash$ be a hash function and $\mathrm{canon}$ a canonical encoding. The \emph{content hash} of $r$ is
\[
\Hash(\mathrm{canon}(\mathsf{type},\mathsf{payload},\mathsf{refs})).
\]
The \emph{record identifier} is $\Hash(\text{content hash},\mathsf{author},\mathsf{sequence})$, where the sequence distinguishes repeated submissions by one author. The \emph{family identifier} is the class of $r$ under the equivalence relation generated by the equivalence claims that are $\IN$ under the current policy.
\end{definition}

The three answer different questions. The content hash supports integrity: has anything changed, and is this the same content? The record identifier names an assertion event: who inscribed it, in what order. The family identifier collects statements proposed to express one proposition; those proposals are themselves contestable claims and never automatic deduplications.

\begin{proposition}[Identity]\label{prop:identity}
Assume $\Hash$ is injective on canonical encodings (the standard idealization of collision resistance, used as a modelling assumption and not proved).
\begin{enumerate}[label=(\alph*)]
\item Two records with identical type, payload and references have the same content hash. If their authors or sequence numbers differ, they have distinct record identifiers.
\item Changing any of type, payload or references changes the content hash and therefore the record identifier.
\item Records with the same content hash lie in the same family.
\end{enumerate}
\end{proposition}

\begin{proof}
(a) and (b) are immediate from injectivity and from the definitions. (c) Identical content is the same proposition, so the equivalence relation of Definition~\ref{def:ids} is taken to contain the relation ``same content hash''.
\end{proof}

Independent confirmations are counted by independence families (Chapter~\ref{ch:evidence}), never by record identifiers. Two authors who submit the same content produce one content hash, two record identifiers and, unless a separate, contestable claim of independence says otherwise, one independence family.

The family identifier is itself an instance of the gluing problem. Equivalence claims are identifications of sign $+1$ between records; a claim that two records express \emph{different} propositions is an identification of sign $-1$. A family is formed from $\IN$ equivalence claims, and an odd cycle mixing the two kinds is exactly a negative cycle in the sense of Theorem~\ref{thm:signed}. The obstruction region of Chapter~\ref{ch:distinction} therefore flags families whose members cannot all be identified as asserted.

\section{Trace completeness}\label{sec:trace}

A contribution attributed to an agent, human or automated, is a claim about an act. It is complete only when the act's causal trace is resolved, and until then it is an attribution and nothing more. An attribution of authorship does not establish that the contribution supports what it claims to support:
\[
\mathrm{ProposedBy}(a,c)\ \not\Rightarrow\ \mathrm{Supports}(a,c).
\]

\begin{definition}[Trace]\label{def:trace}
The \emph{trace} of an act record has six components: the \emph{sources transformed} (evidence items or earlier records the act used), the \emph{warrant applied}, the \emph{scope inferred}, the \emph{earlier records that conditioned the output}, the \emph{independence family}, and the \emph{disposition} that completed the act. A component is \emph{resolved} if it is present and, where it is a reference, points to a record in the ledger. An act is \emph{trace-null}, written $\mathrm{trace}(r)=\bot$, if any component is unresolved.
\end{definition}

The null does not say that the act was worthless. It marks incomplete determination from the position of the ledger: there is a contribution whose derivation cannot yet be inspected. Such a record may be stored, routed for investigation and completed later by anyone able to supply the missing component; it cannot masquerade as completed evidence. The same rule covers text produced by a model with no recorded inputs, an anonymous assertion with no sources, and a bulk import with no methodology, all of which are trace-null for the same reason.

\section{Dispositions and the active set}

Every non-admission has a reason, and the reason is part of the record of what happened.

\begin{definition}[Disposition]\label{def:disposition}
Given a valid ledger $\ledger$ and policy $\policy$, each record $r$ receives a disposition
\begin{align*}
D_\policy(r,\ledger)\in\{{}&\mathrm{ADMITTED},\mathrm{DUPLICATE},\mathrm{MALFORMED},\mathrm{UNSCOPED},\\&\mathrm{QUARANTINED},\mathrm{GATED},\mathrm{REFUSED}\}
\end{align*}
by induction along the reference order: $D_\policy$ depends only on the fields of $r$, on the dispositions of the records it references, on $\policy$, and on gate and capability records the policy names. In particular an argument is $\mathrm{MALFORMED}$ unless well-formed (W1)--(W5) and its fillers are admitted; a claim without scope or modality is $\mathrm{UNSCOPED}$ (Design rule~\ref{rule:scope}); a trace-null act is $\mathrm{QUARANTINED}$; a record whose content hash is shared with a record of smaller record identifier is $\mathrm{DUPLICATE}$ (a canonical tie-break, so that the disposition does not depend on arrival order; on an anchored surface the pair of ledger time and identifier is used, Chapter~\ref{ch:gates}). The \emph{active set} is
\[
\Act(\ledger,\policy)=\{r\in\ledger: D_\policy(r,\ledger)=\mathrm{ADMITTED}\}.
\]
\end{definition}

Every non-admission must state the governing rule, the failing field or argumentative requirement, the evidence for the determination, the responsible decision mechanism and the repair or appeal path. A gated or quarantined record keeps its identity and can become admissible later without resubmission. The argument set $\mathcal A$ of Chapter~\ref{ch:objection} is the set of arguments in $\Act(\ledger,\policy)$, and the labeling is computed on it and on nothing else. This is what ``activation'' means formally.

\begin{proposition}[Inert records do not influence labels]\label{prop:inert}
Let $\ledger,\ledger'$ be valid ledgers with $\Act(\ledger,\policy)=\Act(\ledger',\policy)$. Then the labelings coincide. In particular adding quarantined, gated, duplicate, malformed or unscoped records changes no label.
\end{proposition}

\begin{proof}
The labeling is a function of the finite set of active arguments and of $\policy$ (Definition~\ref{def:grounded}).
\end{proof}

\section{Validity and merging}

\begin{definition}[Valid ledger]\label{def:valid}
A finite set $\ledger$ of records is \emph{valid} if (i) it is \emph{closed}: every reference of every member is a member; (ii) identifiers agree with contents, so that records with equal identifiers are equal; and (iii) there is a linear order of $\ledger$ in which every record follows all the records it references.
\end{definition}

Condition (iii) says that the reference graph is acyclic and is checkable in linear time. Under the hash idealization it is also unavoidable in practice, since a record cannot cite an identifier that depends on its own content; but validity is defined by (iii), and no proof below relies on hashing to obtain acyclicity.

\begin{theorem}[Validity is preserved by appending and merging]\label{thm:valid}
\begin{enumerate}[label=(\alph*)]
\item If $\ledger$ is valid and $r\notin\ledger$ is a record whose references all lie in $\ledger$, then $\ledger\cup\{r\}$ is valid.
\item If $\ledger_1,\ledger_2$ are valid and records with the same identifier in both are equal, then $\ledger_1\cup\ledger_2$ is valid.
\item (\emph{Prefix stability.}) If $\ledger\subseteq\ledger'$ are valid, every record of $\ledger$ has the same content, references and ancestors in $\ledger'$ as in $\ledger$.
\end{enumerate}
\end{theorem}

\begin{proof}
(a) Place $r$ last in a valid order. (c) Identifiers determine records, and by closure every ancestor of a record already lies in $\ledger$.

(b) Closure and identifier agreement are inherited from the two ledgers. For acyclicity, suppose the union's reference graph contained a cycle $r_1\to r_2\to\dots\to r_m\to r_1$, and say $r_1\in\ledger_1$. Every edge leaving a record is determined by that record's references, and $\ledger_1$ is closed, so all the ancestors of $r_1$, hence every $r_i$ on the cycle, lie in $\ledger_1$, with their edges. The cycle lies inside $\ledger_1$, contradicting (iii) there. Any linear extension of the acyclic union graph is the required order.
\end{proof}
